% ----------------------------------------------------------
% Metodologia
% 
% ----------------------------------------------------------

\chapter[Pesquisa/Amostragem]{Pesquisa/Amostragem}

A metodologia aplicada neste trabalho consiste no desenvolvimento de um sistema computacional voltado ao pré-processamento assistido de matrizes de dados metabolômicos de alta dimensionalidade, adquiridas a partir da LC-MS. O escopo do projeto compreende o conjunto de etapas que transformam a matriz semiprocessada em uma matriz limpa, normalizada e estatisticamente robusta, própria para a entrada nas análises estatísticas e biológicas posteriores.
 
\section{Classificação da Pesquisa}
A pesquisa apresenta abordagem predominantemente quantitativa, uma vez que trabalha com dados numéricos gerados por equipamentos analíticos e utiliza métricas estatísticas para avaliar a qualidade do pré-processamento em cada etapa do pipeline. O desempenho dos modelos de aprendizado de máquina também é avaliado por critérios objetivos. 

No que diz respeito aos objetivos, o trabalho assume natureza exploratória e descritiva. É exploratório porque investiga combinações ainda pouco padronizadas entre técnicas quimiométricas e algoritmos de apoio à decisão voltados à assistência do pré-processamento em metabolômica não direcionada. Já a dimensão descritiva aparece no detalhamento sistemático das etapas do fluxo proposto, das decisões algorítmicas tomadas pelo algoritmo e dos critérios usados para validar os resultados.
 
Do ponto de vista dos procedimentos técnicos, o estudo combina duas modalidades. A pesquisa bibliográfica deu base ao referencial teórico apresentado, reunindo as principais abordagens de pré-processamento, normalização, controle de qualidade e análise estatística usadas em metabolômica baseada em LC-MS. A pesquisa experimental envolve o desenvolvimento da ferramenta e a execução de testes com matrizes reais e simuladas, com o objetivo de avaliar como os algoritmos implementados se comportam em diferentes cenários.

\section{Fluxo de Trabalho da Ferramenta}
 
O fluxo de trabalho da ferramenta foi desenhado a partir das etapas críticas de pré-processamento de dados metabolômicos descritas no referencial teórico, com a inserção de pontos de decisão assistidos conduzidos pelo algoritmo de diagnóstico e auxílio à decisão. O processo se inicia com a importação da matriz de dados em formato CSV, obtida a partir da conversão dos arquivos de saída do espectrômetro pelo Progenesis QI, responsável neste fluxo pelo processamento primário (alinhamento, detecção de picos e deconvolução) e pela exportação para um formato tabular interoperável. A matriz contém as \textit{features} metabolômicas organizadas por tempo de retenção, razão massa/carga ($m/z$) e intensidade de sinal por amostra. Após a validação estrutural do arquivo, é executada uma análise exploratória inicial responsável por diagnosticar a qualidade do conjunto, identificando a proporção de valores ausentes, a presença de \textit{outliers}, a variabilidade entre amostras de controle de qualidade e eventuais indícios de deriva instrumental.
 
Com base nesse diagnóstico, o algoritmo sugere o método de correção da variação técnica e a divisão em lotes, deixando a decisão a critério dos pesquisadores. Em seguida, é aplicada a filtragem de \textit{features}, na qual variáveis com mais de 80\% de ausência em todos os grupos ou com RSD elevado nas amostras de controle de qualidade são removidas, em consonância com as boas práticas consolidadas na literatura.
 
\section{Conjunto de Dados e Amostragem}
 
Para a validação experimental da ferramenta, são utilizados conjuntos de dados metabolômicos provenientes de experimentos reais de LC-MS, obtidos em colaboração com o laboratório de metabolômica, complementados, quando necessário, por matrizes sintéticas geradas para reproduzir cenários específicos de degradação dos dados.
  
A estratégia de validação combina, portanto, dados reais --- que asseguram a aderência da ferramenta às condições de uso em laboratório --- e dados sintéticos controlados, que permitem avaliar o comportamento do algoritmo de auxílio à decisão em situações limítrofes que dificilmente seriam encontradas em um único experimento real, ampliando assim a confiabilidade dos testes realizados sobre o pipeline de pré-processamento.